\section{Go Compiler Features}
\label{sec:go-compiler-features}

The Go compiler is the only Emi target that generates a full server, not
just a client for one --- so this section tours what it adds on top of
the shared DTO/action codegen already covered above: database-backed
entities, typed relations with real migratable \href{https://gorm.io}{gorm}
associations, and the create/update actions generated for every entity.
Each part below shows a small sample of the Go code the compiler actually
emits.

\subsection{Entities}
\label{sec:go-entities}

An entity is a database-backed structure: its fields become both a Go
struct and database columns, migrated with \href{https://gorm.io}{gorm}.
Declared in a module's \texttt{entities:} list, right next to
\texttt{dtos:}/\texttt{actions:}:

\begin{lstlisting}
entities:
  - name: entity1
    table: entity1_table
    description: Exercises every supported field type.
    fields:
      - name: title
        type: string
\end{lstlisting}

\texttt{name} becomes the Go struct name (always upper-cased with
\texttt{Entity} appended --- \texttt{Entity1Entity}); \texttt{table}
overrides gorm's default pluralized/snake\_cased table name.

\subsubsection{Two fields every entity gets for free}
\begin{lstlisting}
type Entity1Entity struct {
  Id    int64  `gorm:"primaryKey;autoIncrement" json:"-"`
  UniqueId string `gorm:"type:uuid;default:gen_random_uuid();unique"
                    json:"uniqueId"`
  Title string `json:"title"`
}
\end{lstlisting}
\textbf{\texttt{Id}} is the real gorm primary key --- a plain
auto-incrementing integer, never exposed over JSON/CLI, so joins and
foreign keys can use a small integer instead of a UUID (smaller indexes,
faster joins). \textbf{\texttt{UniqueId}} is the public identifier every
API/CLI surface actually works with: a UUID generated natively by
Postgres itself (\texttt{gen\_random\_uuid()}, built into Postgres 13+)
via a column default, not application code. Both are prepended once, in
\texttt{EntityDefaultFields} (\texttt{lib/golang/go-entity-default-\allowbreak
fields.go}), before the struct generator, the update-input builder, and
the gorm-tag pass all run --- every consumer sees them as ordinary
declared fields.

\end{multicols}

\subsubsection{Field types on gorm}
\label{sec:go-gorm-types}
{\small
\begin{tabular}{@{}p{0.22\linewidth}p{0.74\linewidth}@{}}
\toprule
\textbf{Field type} & \textbf{Gorm tag} \\
\midrule
scalar (\texttt{string}, \texttt{int}, \ldots) & none --- gorm maps it directly \\
\texttt{enum}/\texttt{enum?} & none --- already a plain \texttt{string}/\texttt{Nullable[string]} \\
\texttt{object}/\texttt{object?} & \texttt{embedded} (inlined as columns) \\
\texttt{map}, \texttt{slice} (non-nullable) & \texttt{serializer:json} \\
\texttt{map?}, \texttt{slice?} & none --- \texttt{Nullable[T]} already implements \texttt{Value()}/\texttt{Scan()} \\
\texttt{complex} & none --- the type implements gorm's interfaces itself \\
\texttt{array}/\texttt{collection}/\texttt{one} & see relations, \S\ref{sec:go-relations} \\
\bottomrule
\end{tabular}
}

\begin{multicols}{2}
\raggedcolumns
An explicit \texttt{tags: \{ gorm: ... \}} on a field always wins over
these defaults. A \texttt{complex} field used on an entity must implement
\texttt{driver.Valuer}/\texttt{sql.Scanner} itself --- the generator has
no way to guess how to persist a hand-written type.

Every entity file ends with a small, separate template
(\texttt{go-entity.tpl}) whose output is appended \emph{after} the
generated struct --- currently a \texttt{TableName()} override wired to
\texttt{table:}, and the place to extend if a project wants something
appended to every generated entity file.

\subsubsection{Browsing: JSON Logic filters}
An entity with browsing enabled also gets a generated
\texttt{\{Entity\}BrowseFn} and a \texttt{GET /entity1/browse} action,
exposing five query params: \texttt{filter}, \texttt{sort},
\texttt{startIndex}, \texttt{itemsPerPage}, \texttt{cursor}.
\texttt{filter} is a \href{https://jsonlogic.com}{JSON Logic} expression,
transpiled straight to SQL by \texttt{emigorm.ApplyQueryFilter} --- never
string-concatenated, and never a filter language Emi invented itself:
\begin{lstlisting}
GET /entity1/browse?filter={"==":[{"var":"title"},"hello"]}
GET /entity1/browse?filter={"and":[
  {">":[{"var":"age"},18]},
  {"contains":[{"var":"title"},"draft"]}
]}
\end{lstlisting}
Every standard JSON Logic operator works (\texttt{==}, \texttt{!=},
\texttt{>}, \texttt{<}, \texttt{and}, \texttt{or}, \texttt{in}, \ldots)
plus one Emi adds itself: \texttt{contains} --- a case-insensitive
substring match (\texttt{::text ILIKE '\%...\%'}), the one operator
plain JSON Logic has no native equivalent for.

Sort and paging are the plain, non-JSON-Logic half:
\texttt{sort=createdAt desc} (empty defaults to \texttt{id asc}, needed
for keyset paging to stay correct), and either offset paging
(\texttt{startIndex}/\texttt{itemsPerPage}) or cursor paging
(\texttt{cursor}, echoing the response's own cursor field back on the
next call --- \texttt{cursor=id(42)}).

\begin{pullquote}
\texttt{filter} is the only \emph{user}-controlled half of a Browse
query. A handler's own tenant/workspace isolation goes through the
separate \texttt{emigorm.ApplyQueryScope} (a raw SQL fragment the
handler itself writes, applied unconditionally) --- \texttt{filter} has
no way to see or override it.
\end{pullquote}

\subsection{Writing a complex type}
\label{sec:go-complex}
A \texttt{complex: Money} field (\S\ref{sec:complex-types}) needs a real
Go type at the declared \texttt{location}/\texttt{namespace} --- nothing
more, to just carry the value through a DTO or action. Persisting it on
an \texttt{entity} is a separate, opt-in step: gorm has no idea how to
store an arbitrary hand-written struct, so the type has to say so itself.

\subsubsection{Nullability: no wrapper, no pointer}
A \texttt{complex} field is the one place Go's nullable-field convention
doesn't apply. Every other nullable field type gets wrapped in
\texttt{emigo.Nullable[T]} (\S\ref{sec:go-gorm-types}); a complex field
never does, whether the field is declared \texttt{complex} or (the
JS-only variant) \texttt{complex?} --- both compile to the exact same
plain \texttt{Money} value in Go, because a complex value has no
wire-level nullability of its own as far as the Go compiler is concerned.
\texttt{complex?} exists purely so the JS/TS compiler can tell ``always
instantiate'' apart from ``allow \texttt{undefined}/\texttt{null} on the
setter''; every other backend, Go included, treats the two identically.

The one place a complex field \emph{does} become a pointer is inside the
generated \texttt{\{Entity\}UpdateInput} struct (\S\ref{sec:go-actions}):
since there's no \texttt{Nullable[T]} counterpart for a hand-written
type, the update-input builder falls back to a plain \texttt{*Money},
where \texttt{nil} means ``untouched'' --- the same ``was this field
mentioned at all'' signal every other field gets from
\texttt{Nullable[T]}'s own \texttt{IsSet()}, just spelled with a raw
pointer instead of a generic wrapper.

\subsubsection{Persisting to one column}
Implementing \texttt{database/sql/driver.Valuer} and \texttt{sql.Scanner}
stores the whole value as a single column:
\begin{lstlisting}
type Money struct {
  Cents int64
}

func (m Money) Value() (driver.Value, error) {
  return m.Cents, nil
}

func (m *Money) Scan(value interface{}) error {
  switch v := value.(type) {
  case int64:
    m.Cents = v
  case []byte:
    cents, err := strconv.ParseInt(string(v), 10, 64)
    if err != nil {
      return err
    }
    m.Cents = cents
  }
  return nil
}
\end{lstlisting}
A real complex type is expected to implement whatever interfaces its own
storage needs --- \texttt{MarshalText}/\texttt{UnmarshalText} too, if it
should also cast cleanly from a CLI string argument, the same way
fireback's own complex types (e.g.\ \texttt{XDate}) do.

\subsubsection{Persisting across multiple columns}
Valuer/Scanner is a single-column round trip; not every complex type fits
one column. Skip both interfaces and gorm falls back to walking the
struct like it does for an \texttt{object}/\texttt{object?} field
(\S\ref{sec:go-gorm-types}) --- each exported field becomes its own column:
\begin{lstlisting}
type Address struct {
  Street string `gorm:"column:addr_street"`
  City   string `gorm:"column:addr_city"`
  Zip    string `gorm:"column:addr_zip"`
}
\end{lstlisting}
No \texttt{Value()}/\texttt{Scan()} at all here --- the struct is plain
data, and gorm's own reflection spreads it across \texttt{addr\_street},
\texttt{addr\_city}, \texttt{addr\_zip} exactly as it would for any
embedded struct. Pick single-column when the type has one natural
serialized form (money as cents, a date as text); pick multi-column when
it's naturally shaped as several fields side by side.
